\setchoices{label=(\Alph*)}

\begin{question}[index=1]
设函数 $f(x)$ 在区间 $(-1,1)$ 内有定义，且在点 $x=0$ 处连续，则以下结论：
\par\circlednumber{1} 当 $\displaystyle\lim_{x\to0}\frac{f(x)}{\sqrt[3]x}=0$ 时，$f(x)$ 在点 $x=0$ 处可导；
\par\circlednumber{2} 当 $\displaystyle\lim_{x\to0}\frac{f(x)}{x^2}=0$ 时，$f(x)$ 在点 $x=0$ 处可导；
\par\circlednumber{3} 当 $f(x)$ 在点 $x=0$ 处可导时，$\displaystyle\lim_{x\to0}\frac{f(x)}{\sqrt[3]x}=0$.
\par 所有正确结论的序号为
\begin{choices}[columns=4]
\item ①
\item ②
\item ②③
\item ①②
\end{choices}
\end{question}
\QuestionBreak

\begin{question}[index=2]
方程 $\eu^{z+1}+xy+yz+zx=0$ 在点 $(1,1,-1)$ 的某个邻域内必能确定的隐函数以及该隐函数确定的曲面在该点的切平面方程为
\begin{choices}[columns=2]
\item $z=z(x,y),\ z=-1$
\item $x=x(y,z)$ 与 $y=y(z,x),\ x+y=2$
\item $z=z(x,y),\ z=1$
\item $x=x(y,z)$ 与 $y=y(z,x),\ x-y=2$
\end{choices}
\end{question}
\QuestionBreak

\begin{question}[index=3]
设常数 $\alpha>0$，正项级数 $\displaystyle\sum_{n=1}^\infty a_n$ 收敛，则级数 $\displaystyle\sum_{n=1}^\infty(-1)^n\frac{\sqrt{a_{n+2}}}{\sqrt{n^2+1+\cos\alpha}}$
\begin{choices}[columns=2]
\item 发散
\item 条件收敛
\item 绝对收敛
\item 敛散性与 $\alpha$ 有关
\end{choices}
\end{question}
\QuestionBreak

\begin{question}[index=4]
设函数 $f(u)$ 连续，记 $\displaystyle I=\iint_Df(x+y)\dd x\dd y$，其中 $D=\{(x,y)\mid x\ge0,y\ge0,x+y\le a\}$（$a>0$），则 $I=$
\begin{choices}[columns=2]
\item $\displaystyle\int_0^axf(x)\dd x$
\item $\displaystyle\int_a^0xf(x)\dd x$
\item $\displaystyle\int_0^a(a-x)f(x)\dd x$
\item $\displaystyle\int_a^0(a-x)f(x)\dd x$
\end{choices}
\end{question}
\QuestionBreak

\begin{question}[index=5]
已知 $\mat{1&1&1\\1&2&1\\1&1&1}=\mat{1&1\\1&2\\1&1}\vect A$，则 $\vect A=$
\begin{choices}[columns=2]
\item $\mat{1&0&1\\0&0&1}$
\item $\mat{0&1&0\\1&0&1}$
\item $\mat{1&1&0\\0&0&1}$
\item $\mat{1&0&1\\0&1&0}$
\end{choices}
\end{question}
\QuestionBreak

\begin{question}[index=6]
设 $\lambda_1,\lambda_2$ 是 $3$ 阶矩阵 $\vect A$ 的两个不同特征值，$\vect\alpha_1,\vect\alpha_2$ 是 $\vect A$ 的属于 $\lambda_1$ 的线性无关的特征向量，$\vect\alpha_3$ 是 $\vect A$ 的属于 $\lambda_2$ 的特征向量，则 $\vect\alpha_1+\vect A\vect\alpha_3$，$\vect A(\vect\alpha_2-\vect\alpha_3)$，$\vect A\vect\alpha_1+\vect\alpha_3$ 线性相关的充要条件是
\begin{choices}[columns=2]
\item $\lambda_1=0$ 或 $\lambda_1\lambda_2=1$
\item $\lambda_2=0$ 或 $\lambda_1\lambda_2=1$
\item $\lambda_1\ne0$ 且 $\lambda_1\lambda_2\ne1$
\item $\lambda_2\ne0$ 且 $\lambda_1\lambda_2\ne1$
\end{choices}
\end{question}
\QuestionBreak

\begin{question}[index=7]
若 $\vect A$ 为 $n$ 阶矩阵，且满足 $\vect A^2-\vect A-2\vect E=\vect O$，则
\begin{choices}[columns=2]
\item $\vect A$ 为对称矩阵
\item $\vect A=2\vect E$
\item $\vect A$ 可相似对角化
\item $\vect A$ 的特征值必有 $2$
\end{choices}
\end{question}
\QuestionBreak

\begin{question}[index=8]
设 $(X,Y)$ 服从二维正态分布 $N(0,0;1,1;\rho)$，令 $U=X+Y$，$V=X-Y$，则 $\operatorname{Cov}(U,V)=$
\begin{choices}[columns=4]
\item $2$
\item $0$
\item $2+2\rho$
\item $2-2\rho$
\end{choices}
\end{question}
\QuestionBreak

\begin{question}[index=9]
设 $X_1,X_2$ 是取自正态总体 $X\dist N(1,1)$ 的简单随机样本，则 $\dfrac{X_1-1}{|1-X_2|}$ 服从
\begin{choices}[columns=4]
\item $N(1,1)$
\item $\chi^2(1)$
\item $F(1,1)$
\item $t(1)$
\end{choices}
\end{question}
\QuestionBreak

\begin{question}[index=10]
设总体 $X\dist N(\mu,\sigma^2)$，其中 $\mu,\sigma$ 均未知，对于假设检验问题：$H_0:\sigma^2\le8$；$H_1:\sigma^2>8$，给定显著性水平 $\alpha=0.05$，对于一个容量为 $25$ 的样本，计算得样本方差 $s^2=14$，已知 $\chi_{0.05}^2(24)=36.415$，则检验结果为
\begin{choices}[columns=2]
\item 接受 $H_0$，可能会犯第一类错误
\item 接受 $H_0$，可能会犯第二类错误
\item 拒绝 $H_0$，可能会犯第一类错误
\item 拒绝 $H_0$，可能会犯第二类错误
\end{choices}
\end{question}
